\documentclass[11pt]{book}

\usepackage[english]{babel}
\usepackage{caption}
\usepackage{cite}
\usepackage{color}
\usepackage[T1]{fontenc}
\usepackage{graphicx}
\usepackage[utf8]{inputenc}
\usepackage[hidelinks]{hyperref}
\usepackage{listings}
\usepackage{listings-rust}
\usepackage{newclude}
\usepackage{pdfpages}
\usepackage{todonotes}

\setlength{\textwidth}{146.8mm} % = 210mm - 37mm - 26.2mm
\setlength{\oddsidemargin}{11.6mm} % 37mm - 1in (from hoffset)
\setlength{\evensidemargin}{0.8mm} % = 26.2mm - 1in (from hoffset)
\setlength{\topmargin}{-2.2mm} % = 0mm -1in + 23.2mm
\setlength{\textheight}{221.9mm} % = 297mm -29.5mm -31.6mm - 14mm (12 to accomodate footline with pagenumber)
\setlength{\headheight}{14pt}
\setlength{\parindent}{0pt}

\graphicspath{{../diagram}{../graphs}}

\usepackage{setspace}
\setstretch{1.1}

\newcommand{\marios}[1]{\noindent{{\bf \fbox{MK:} {\textcolor{green}{\it#1}}}}}

\makeatletter
\setlength{\@fptop}{0pt}
\makeatother

\usepackage{fancyhdr}
\renewcommand{\sectionmark}[1]{\markright{\thesection\ #1}}
\pagestyle{fancy}
\fancyhf{}
\renewcommand{\headrulewidth}{0.4pt}
\renewcommand{\footrulewidth}{0pt}
\fancyhead[OR]{\bfseries \nouppercase{\rightmark}}
\fancyhead[EL]{\bfseries \nouppercase{\leftmark}}
\fancyfoot[EL,OR]{\thepage}
\fancypagestyle{plain}{
  \fancyhf{}
  \renewcommand{\headrulewidth}{0pt}
  \renewcommand{\footrulewidth}{0pt}
  \fancyfoot[EL,OR]{\thepage}}
\fancypagestyle{addpagenumbersforpdfimports}{
  \fancyhead{}
  \renewcommand{\headrulewidth}{0pt}
  \fancyfoot{}
  \fancyfoot[RO,LE]{\thepage}
}

\usepackage{tikz}
\usepackage[explicit]{titlesec}
\newcommand*\chapterlabel{}
% \renewcommand{\thechapter}{\Roman{chapter}}
\titleformat{\chapter}[display]  % type (section,chapter,etc...) to vary,  shape (eg display-type)
{\normalfont\bfseries\Huge} % format of the chapter
{\gdef\chapterlabel{\thechapter\ }}     % the label
{0pt} % separation between label and chapter-title
{\begin{tikzpicture}[remember picture,overlay]
    \node[yshift=-8cm] at (current page.north west)
    {\begin{tikzpicture}[remember picture, overlay]
        \draw[fill=black] (0,0) rectangle(35.5mm,15mm);
        \node[anchor=north east,yshift=-7.2cm,xshift=34mm,minimum height=30mm,inner sep=0mm] at (current page.north west)
        {\parbox[top][30mm][t]{15mm}{\raggedleft $\phantom{\textrm{l}}$\color{white}\chapterlabel}};  %the black l is just to get better base-line alingement
        \node[anchor=north west,yshift=-7.2cm,xshift=37mm,text width=\textwidth,minimum height=30mm,inner sep=0mm] at (current page.north west)
        {\parbox[top][30mm][t]{\textwidth}{\color{black}#1}};
      \end{tikzpicture}
    };
  \end{tikzpicture}
  \gdef\chapterlabel{}
} % code before the title body

\titlespacing*{\chapter}{0pt}{50pt}{30pt}
\titlespacing*{\section}{0pt}{13.2pt}{*0}  % 13.2pt is line spacing for a text with 11pt font size
\titlespacing*{\subsection}{0pt}{13.2pt}{*0}
\titlespacing*{\subsubsection}{0pt}{13.2pt}{*0}

\newcounter{myparts}
\newcommand*\partlabel{}
\titleformat{\part}[display]  % type (section,chapter,etc...) to vary,  shape (eg display-type)
{\normalfont\bfseries\Huge} % format of the part
{\gdef\partlabel{\thepart\ }}     % the label
{0pt} % separation between label and part-title
{\setlength{\unitlength}{20mm}
  \addtocounter{myparts}{1}
  \begin{tikzpicture}[remember picture,overlay]
    \node[anchor=north west,xshift=-65mm,yshift=-6.9cm-\value{myparts}*20mm] at (current page.north east) % for unknown reasons: 3mm missing -> 65 instead of 62
    {\begin{tikzpicture}[remember picture, overlay]
        \draw[fill=black] (0,0) rectangle(62mm,20mm);   % -\value{myparts}\unitlength
        \node[anchor=north west,yshift=-6.1cm-\value{myparts}*20mm,xshift=-60.5mm,minimum height=30mm,inner sep=0mm] at (current page.north east)
        {\parbox[top][30mm][t]{55mm}{\raggedright \color{white}Part \partlabel $\phantom{\textrm{l}}$}};  %the phantom l is just to get better base-line alingement
        \node[anchor=north east,yshift=-6.1cm-\value{myparts}*20mm,xshift=-63.5mm,text width=\textwidth,minimum height=30mm,inner sep=0mm] at (current page.north east)
        {\parbox[top][30mm][t]{\textwidth}{\raggedleft \color{black}#1}};
      \end{tikzpicture}
    };
  \end{tikzpicture}
  \gdef\partlabel{}
} % code before the title body

\usepackage{etoolbox}
\let\bbordermatrix\bordermatrix
\patchcmd{\bbordermatrix}{8.75}{4.75}{}{}
\patchcmd{\bbordermatrix}{\left(}{\left[}{}{}
    \patchcmd{\bbordermatrix}{\right)}{\right]}{}{}
\usepackage{amsmath}
% Fix the problem with delimiter size caused by fourier and amsmath packages.
\makeatletter
\def\resetMathstrut@{%
  \setbox\z@\hbox{%
    \mathchardef\@tempa\mathcode`\(\relax
    \def\@tempb##1"##2##3{\the\textfont"##3\char"}%
    \expandafter\@tempb\meaning\@tempa \relax
  }%
  \ht\Mathstrutbox@1.2\ht\z@ \dp\Mathstrutbox@1.2\dp\z@
}
\makeatother

\lstset{
  backgroundcolor=\color{white},
  basicstyle=\footnotesize,
  breaklines=false,
  deletekeywords= {stack},
  language=Rust,
  frame=single,
  keywordstyle=\color{green},
}

\author{Ogier Bouvier}
\date{\today}
\title{A zero-copy key-value store in Rust}

\begin{document}

\begin{titlepage}
  \newcommand{\HRule}{\rule{\linewidth}{0.5mm}} % Defines a new command for the horizontal lines, change thickness here
  \center

  \LARGE \textsc{École Polytechnique Fédérale de Lausanne}
  \vspace{1cm}

  \LARGE Master Thesis
  \vspace{2cm}

  {
    \HRule \\[0.5cm]
    \Huge A zero copy key-value store in Rust}
  \HRule \\[1cm]

  \begin{minipage}[t]{0.4\textwidth}
    \large
    {\large \today}
  \end{minipage}\\[1cm]

  \vspace{2cm}

  \begin{tabular}{c c}
    \begin{minipage}[t]{0.4\textwidth}
      \begin{flushleft} \large
        \textbf{Author} \\
        Ogier \textsc{Bouvier}\\
        \textit{ogier.bouvier@epfl.ch}
      \end{flushleft}
    \end{minipage}

    \begin{minipage}[t]{0.4\textwidth}
      \begin{flushright} \large
        \textbf{Supervisors} \\
        Prof. Edouard \textsc{Bugnion}\\
        DCSL | EPFL \\
        \textit{edouard.bugnion@epfl.ch} \\
        Marios \textsc{Kogias}\\
        DCSL | EPFL \\
        \textit{marios.kogias@epfl.ch}
      \end{flushright}
    \end{minipage}
  \end{tabular}

  \begin{center}
    \begin{figure}[h]
      \includegraphics[width=0.71\textwidth]{../diagram/epfl.pdf}
    \end{figure}
  \end{center}

\end{titlepage}

\chapter*{Acknowledgements}
\chapter*{Abstract}
\marios{You need a short abstract that describes the problem, hints the solutions and gives some basic performance numbers.}

\pagenumbering{gobble}
\tableofcontents

\listoftodos{}

\pagenumbering{arabic}
\setcounter{page}{1}

\chapter{Introduction}

\include*{intro}

\chapter{Background}

\include*{background}

\chapter{Design}

\include*{design}

\chapter{Implementation}

\include*{impl}

\chapter{Evaluation}

\include*{eval}

\chapter{Further work}

\let\clearpage\relax
\include*{work}

\chapter{Conclusion}

We hope to have provided an easy to use and clear API to build upon
our system. We also think that we have proved that Rust is a valid
language to build kernel bypass networking in. \todo{speak about the
  actual performance}

\pagenumbering{gobble}
\bibliographystyle{plain}
\bibliography{master}{}

\end{document}
%%% Local Variables:
%%% mode: latex
%%% TeX-master: t
%%% End:
